%! The Inverse and Determinant of a Matrix — Unit 4, Lesson 4.11 — Lesson Plan
\documentclass[10pt]{article}
\usepackage{precalculus-article}
\usepackage{precalculus-boxes}
\usepackage{graphicx}
\graphicspath{{images/}}
\newcommand{\TallMath}[1]{$\displaystyle #1\rule[-1.4em]{0pt}{3.2em}$}
\newcommand{\CourseName}{AP Precalculus}
\newcommand{\SchoolYear}{2026--2027}
\newcommand{\MeetingLength}{55 minutes}
\newcommand{\UnitNumberName}{Unit 4: Functions Involving Parameters, Vectors, and Matrices \quad}
\newcommand{\LessonNumberName}{Lesson 4.11: The Inverse and Determinant of a Matrix}

\begin{document}
\begin{center}
    {\Large\bfseries \CourseName: \SchoolYear \\
    {\small\bfseries \UnitNumberName \LessonNumberName}}
\end{center}

\begin{tcolorbox}[colback=sky, colframe=navy, boxrule=0.9pt, arc=2mm,
    left=2mm, right=2mm, top=1.5mm, bottom=1.5mm]
    \textbf{Primary Objective:} Students will compute the determinant of a $2\times2$ matrix,
    use it to determine invertibility, and find the inverse of an invertible $2\times2$ matrix
    using the algebraic formula. (Big Idea: Functions)
\end{tcolorbox}

\renewcommand{\arraystretch}{1.6}
\begin{skillbox}[Priority Ideas \& Skills]{goldbox}
\begin{minipage}[t]{0.48\linewidth}\vspace{0pt}
    \textbf{AP Skills:}
    \begin{itemize}[leftmargin=1.2em, itemsep=2pt]
        \item \textbf{1.B} --- Identify mathematical information from graphical, numerical,
              analytical, and verbal representations.
        \item \textbf{3.B} --- Apply appropriate mathematical procedures for a given situation.
    \end{itemize}
    \smallskip
    \textbf{Learning Objectives:} 4.11.A, 4.11.B\\
    \textbf{Essential Knowledge:} 4.11.A.1--4, 4.11.B.1--3
\end{minipage}\hfill
\begin{minipage}[t]{0.48\linewidth}\vspace{0pt}
    \textbf{Key Understandings:}
    \begin{itemize}[leftmargin=1.2em, itemsep=2pt]
        \item The identity matrix $I$ satisfies $A \cdot I = I \cdot A = A$.
        \item $A^{-1}$ exists if and only if $\det(A) \ne 0$.
        \item $\det(A) = ad - bc$ for $A = \begin{bmatrix}a&b\\c&d\end{bmatrix}$.
        \item $A^{-1} = \dfrac{1}{\det(A)}\begin{bmatrix}d&-b\\-c&a\end{bmatrix}$.
        \item $|\det(A)|$ equals the area of the parallelogram spanned by the columns of $A$.
    \end{itemize}
\end{minipage}
\end{skillbox}

\begin{skillbox}[{Vocabulary, Concepts, \& Theorems}]{greenbox}
\renewcommand{\arraystretch}{1.5}
\begin{tabularx}{\linewidth}{@{} l X @{}}
    \textbf{identity matrix} & A square matrix $I$ with 1s on the main diagonal and 0s elsewhere;
        $I_2 = \begin{bmatrix}1&0\\0&1\end{bmatrix}$. \\
    \textbf{inverse matrix} & $A^{-1}$: the unique matrix satisfying $A \cdot A^{-1} = A^{-1} \cdot A = I$. \\
    \textbf{determinant} & A scalar $\det(A) = ad - bc$ encoding the scaling and orientation of $A$. \\
    \textbf{invertible} & A matrix is \textit{invertible} when $\det(A) \ne 0$; \textit{singular} when $\det(A) = 0$. \\
\end{tabularx}
\end{skillbox}

\begin{skillbox}[Activate Prior Knowledge \& Spiral Review]{sky}
    \begin{tabularx}{\linewidth}{@{}X c@{}}
        \begin{minipage}[t]{0.58\linewidth}\vspace{0pt}
            \textbf{Students should already know:}
            \begin{itemize}[leftmargin=1.2em, itemsep=2pt]
                \item Matrix multiplication (Lesson 4.10): row $\times$ column dot products.
                \item Multiplicative inverses of real numbers: $a \cdot a^{-1} = 1$.
                \item Basic $2\times2$ arithmetic and signed arithmetic.
            \end{itemize}
        \end{minipage}
        &
        \begin{minipage}[t]{0.36\linewidth}\vspace{0pt}\centering
            \textit{See attached warm-up.}
        \end{minipage}
    \end{tabularx}
\end{skillbox}

\begin{skillbox}[Hook]{sky}
    \textbf{Opening question:} ``What number times 6 gives 1?'' (Answer: $\tfrac{1}{6}$.)
    That is the \emph{multiplicative inverse} of 6. Now ask: \textbf{``Can a matrix have an
    inverse --- a matrix $M^{-1}$ such that $M \cdot M^{-1} = I$?''}
    Allow 1--2 minutes of think-pair-share. Reveal: yes, but \emph{only when a certain
    condition is met}. Today we find that condition and derive the formula.
\end{skillbox}

\begin{skillbox}[Lesson]{sky}
\begin{multicols}{2}

\textbf{Step 1 --- The Identity Matrix (EK 4.11.A.1--2).}\\
Display $I_2 = \begin{bmatrix}1&0\\0&1\end{bmatrix}$.
Have students multiply $\begin{bmatrix}3&1\\2&5\end{bmatrix}\begin{bmatrix}1&0\\0&1\end{bmatrix}$
and observe the product equals the original matrix.
Generalize: $A \cdot I = I \cdot A = A$ for any square $A$.

\textbf{Step 2 --- Define the Inverse (EK 4.11.A.3).}\\
We seek $A^{-1}$ so $A \cdot A^{-1} = I$. Ask: ``Why is this useful?''
(Solving matrix equations, like $x = a^{-1} \cdot b$ for real numbers.)
Verify together: $\begin{bmatrix}2&1\\5&3\end{bmatrix}\begin{bmatrix}3&-1\\-5&2\end{bmatrix}=I$.

\columnbreak

\textbf{Step 3 --- The Determinant (EK 4.11.B.1, 4.11.B.3).}\\
Introduce $\det(A) = ad - bc$.
Compute $\det\begin{bmatrix}2&1\\5&3\end{bmatrix} = 6-5 = 1$.
Note: if $\det(A) = 0$, the inverse does not exist.

\textbf{Step 4 --- The Inverse Formula (EK 4.11.A.4).}\\
Present: $A^{-1} = \tfrac{1}{\det(A)}\begin{bmatrix}d&-b\\-c&a\end{bmatrix}$.
Worked example: $A=\begin{bmatrix}3&2\\1&4\end{bmatrix}$, $\det=10$,
$A^{-1}=\tfrac{1}{10}\begin{bmatrix}4&-2\\-1&3\end{bmatrix}$.
Verify $A \cdot A^{-1} = I$ together.

\textbf{Step 5 --- Geometric Interpretation (EK 4.11.B.2).}\\
The columns of $A$ span a parallelogram; $|\det(A)|$ is its area.
Example: $\begin{bmatrix}2&0\\1&3\end{bmatrix}$ has area $|6-0|=6$.
If $\det=0$, the columns are parallel (degenerate, zero area).

\end{multicols}
\end{skillbox}

\begin{skillbox}[Active Monitoring]{redbox}
    \textbf{Circulate and check for:}
    \begin{itemize}[leftmargin=1.2em, itemsep=2pt]
        \item Computing $bc - ad$ instead of $ad - bc$.
        \item Sign errors on $-b$ and $-c$ in the inverse formula.
        \item Forgetting to scale the entire matrix by $\tfrac{1}{\det(A)}$.
        \item Swapping $a$ and $d$ (confusing the swap-and-negate formula).
    \end{itemize}
    \textbf{Cold-call prompts:}
    \begin{itemize}[leftmargin=1.2em, itemsep=2pt]
        \item ``How do you know whether a matrix is invertible \emph{before} computing $A^{-1}$?''
        \item ``If $\det(A) = 0$, what does that mean geometrically?''
    \end{itemize}
\end{skillbox}

\begin{skillbox}[Group Work \& Differentiation]{redbox}
\begin{multicols}{3}
\textbf{Tier R --- Remediate}
\begin{itemize}[leftmargin=1em, itemsep=2pt]
    \item Given $A=\begin{bmatrix}3&2\\1&4\end{bmatrix}$, identify $a,b,c,d$.
    \item Compute $\det(A) = ad - bc$.
    \item Is $A$ invertible? Why or why not?
\end{itemize}

\columnbreak
\textbf{Tier A --- Approaching}
\begin{itemize}[leftmargin=1em, itemsep=2pt]
    \item Find $A^{-1}$ for $A=\begin{bmatrix}3&2\\1&4\end{bmatrix}$.
    \item Verify $A \cdot A^{-1} = I$ by multiplication.
\end{itemize}

\columnbreak
\textbf{Tier E --- Extension}
\begin{itemize}[leftmargin=1em, itemsep=2pt]
    \item Let $B=\begin{bmatrix}3&6\\1&2\end{bmatrix}$.
    \item Compute $\det(B)$. What does $\det(B)=0$ mean geometrically for the column vectors of $B$?
    \item Can $B$ be inverted? Explain why or why not.
\end{itemize}
\end{multicols}
\end{skillbox}

\begin{skillbox}[Individual Work \& Assessment]{redbox}
    \textbf{Exit Ticket (last 8 minutes):}
    \begin{enumerate}[leftmargin=1.5em, itemsep=3pt]
        \item Compute $\det\begin{bmatrix}5&2\\3&1\end{bmatrix}$.
        \item Does this matrix have an inverse? If so, find it.
        \item A $2\times2$ matrix has $\det = 0$. What does this mean geometrically about
              its column vectors?
    \end{enumerate}
    \smallskip
    \textbf{Collect and review:} check for sign errors in the determinant and students who
    attempt to invert a singular matrix. Use results to inform Day 2 pacing.
\end{skillbox}

\begin{skillbox}[Reinforcement \& Extension]{goldbox}
    \textbf{Homework:} Problems covering determinant computation, invertibility checks,
    finding $A^{-1}$, verifying by multiplication, geometric area interpretation, and one
    AP-style multiple-choice item. See \texttt{homework/main.tex}.

    \textbf{Extension:} Explore $3\times3$ determinants via cofactor expansion.

    \textbf{Preview of Lesson 4.12:} Linear transformations --- matrices as functions that
    encode geometric actions (rotation, reflection, scaling). Ask students: ``If $A$ maps
    every point in the plane, what does $A^{-1}$ do?''
\end{skillbox}

\end{document}
